\section{실험 환경}
\label{sec:실험환경}

\subsection{시뮬레이터 및 데이터셋 규모}
본 제안 기법의 성능 검증을 위해, Python 기반의 FANET 라우팅 전용 시뮬레이터인 Lite-GLOBE를 이용해 대규모 평가를 수행하였다. 실험 데이터셋은 총 5개의 훈련용 난수 시드, 14개의 서로 다른 평가용 시나리오, 시나리오당 200회의 에피소드로 구성된다. 평가 기간 동안 총 6가지 알고리즘 변종을 비교하여 약 84,000개의 에피소드 레코드를 획득하고, 이를 집계하여 최종 통계적 신뢰성을 확보하였다.

\subsection{평가 시나리오 구성}
실험 시나리오는 다음과 같은 네 가지 대표 그룹으로 분류된다.
\begin{itemize}
    \item \textbf{일반 이동성 시나리오 (General Unseen Mobility)}: 학습 단계에서 보지 못한 새로운 UAV 기동 궤적 하에서 기본적인 패킷 전송 적응력을 평가한다.
    \item \textbf{라우팅 홀 시나리오 (Routing Void/Hole)}: 목적지 방향으로 포워딩 노드가 존재하지 않아 로컬 미니멈이 발생하는 공간적 장애물이 위치한 환경을 시험한다.
    \item \textbf{예측 단절 시나리오 (Predictive Link Break)}: 상대 위치 상으로는 목적지와 가장 가깝지만, 두 노드의 진행 방향이 반대여서 즉각적으로 연결이 해제되는 고속 기동 상황을 모사한다.
    \item \textbf{노드 확장성 스윕 (Node Scalability Sweep)}: 무인기 대수를 증가시켜 이웃 노드 밀도 변화에 따른 오버헤드와 생존력을 검증한다.
\end{itemize}

\subsection{비교 알고리즘군}
비교 평가를 위해 총 5개의 외부 대표 알고리즘 및 2개의 자체 아블레이션 모델을 구현하였다.
\begin{itemize}
    \item \textbf{GPSR}: 1-hop 물리 좌표만을 사용하는 고전적인 지리적 그리디 라우팅 프로토콜.
    \item \textbf{Predictive Geographic}: 링크 잔여 수명 등 수학적 모빌리티 예측 모형을 결합한 지리적 라우팅 프로토콜.
    \item \textbf{Evo-QGeo}: 미래 링크 상태의 주기적 진화를 고려하는 외부 강화학습 계열 예측형 라우팅 프로토콜.
    \item \textbf{IQMR Q($\lambda$)}: TD 학습 방식 및 다중 홉 보상 갱신을 수행하는 외부 Q-러닝 라우팅 프로토콜.
    \item \textbf{DRAMA}: 에이전트 간 런타임 제어 메시지 교환 및 창발적 통신 기법을 채택한 외부 MARL 라우팅 프로토콜.
    \item \textbf{Phase 8 Geo-Residual KD}: 위험 스위칭이 결여되고 지리적 우선도 Prior에 오프라인 모방 잔차만을 훈련시킨 자체 기본 학생 모델.
\end{itemize}

\subsection{평가 지표}
기본적인 신뢰성 검증을 위해 패킷 전달 성공률(PDR)과 데드라인 내 성공 전송 비율을 산출한다. 또한 전송 지연 성능의 tail-behavior를 확인하기 위해 평균 delay와 더불어 95% 백분위수 지연 시간(Delay$_{95}$)을 비교 분석한다. 시스템 효율을 증명하기 위해 송신 시 사용된 에너지 대리 수치 및 차홉 결정 단계에서 알고리즘이 소비한 입력 및 제어 정보 크기(Byte)를 동시에 리포트한다.
